titulo: Demostración del Teorema Fundamental del Cálculo
tema: Integración
dificultad: intermedia
etiquetas: integral, continuidad, derivada
actualizada: 2026-07-31
---
\begin{teorema}
Sea $f$ integrable en $[a,b]$ y definamos $F(x)=\int_a^x f(t)\,dt$.
Si $f$ es continua en $c \in [a,b]$, entonces $F$ es derivable en $c$ y
\[
  F'(c) = f(c).
\]
\end{teorema}

\begin{demostracion}
Para $h \neq 0$ con $c+h \in [a,b]$,
\[
  \frac{F(c+h)-F(c)}{h} = \frac{1}{h}\int_c^{c+h} f(t)\,dt .
\]
Sea $\varepsilon > 0$. Por continuidad de $f$ en $c$ existe $\delta>0$ tal que
$|t-c|<\delta$ implica $|f(t)-f(c)|<\varepsilon$. Entonces, para $0<|h|<\delta$,
\[
  \left| \frac{1}{h}\int_c^{c+h} f(t)\,dt - f(c) \right|
  = \left| \frac{1}{h}\int_c^{c+h} \bigl(f(t)-f(c)\bigr)\,dt \right|
  \le \varepsilon .
\]
Por lo tanto el límite existe y vale $f(c)$. \qed
\end{demostracion}
